\documentclass[11pt,a4paper]{article}
\usepackage{dokumentacja}

\nazwaprojektu{Sale do nauki}
\podtytulprojektu{Aplikacja, w której student znajduje wolną salę na konkretną godzinę i ją rezerwuje}
\tematprojektu{7. Rezerwacja sal do nauki}
\nazwazespolu{}
\numerzespolu{7}
\grupalab{LB2}
\repozytorium{}
\technologia{HTML, CSS, JavaScript (bez frameworka)}
\wersjadokumentu{0.1 - koncepcja i makiety (K1)}
\dataoddania{25.11.2026}

\osoba{Szymon Nowomiński}{177924}{ekran startowy, nawigacja, lista i szczegóły sali}
\osoba{Bartosz Ogórek}{177941}{formularz z walidacją, moje rezerwacje, wąski ekran}

\hypersetup{pdftitle={Dokumentacja projektu},pdfsubject={Metody opracowywania interfejsów użytkownika}}

\begin{document}

\stronatytulowa

\section*{Metryczka dokumentu}
\addcontentsline{toc}{section}{Metryczka dokumentu}

\subsection*{Dane projektu}
\tabeladanych

\subsection*{Skład zespołu}
\tabelaskladu

\subsection*{Historia wersji}
\begin{tabularx}{\linewidth}{@{}p{1.6cm}p{2.4cm}p{3.6cm}L@{}}
\naglowektabeli \nt{Wersja} & \nt{Data} & \nt{Autor} & \nt{Opis zmian} \\
0.1 & 17.11 & Szymon Nowomiński, Bartosz Ogórek & koncepcja i makiety (punkt kontrolny K1) \\
\bottomrule
\end{tabularx}

\subsection*{Oświadczenie zespołu}
Oświadczamy, że projekt i dokumentacja są wynikiem pracy naszego zespołu oraz że każda osoba potrafi wyjaśnić dowolny fragment kodu.

\newpage


\tableofcontents
\listoffigures
\listoftables
\newpage

\setcounter{section}{1}
\section{Użytkownik i cel}\label{sec:uzytkownik}

Student, który chce pouczyć się w grupie na terenie uczelni. Jego głównym celem jest znaleźć wolną salę na konkretną godzinę i ją zarezerwować.
\setcounter{section}{3}
\section{Ekrany i nawigacja}\label{sec:nawigacja}

\subsection{Lista ekranów}
\begin{table}[H]\centering
\caption{Ekrany aplikacji}\label{tab:ekrany}
\begin{tabularx}{\linewidth}{@{}p{0.6cm}p{3.6cm}Lp{3.4cm}@{}}
\naglowektabeli \nt{\#} & \nt{Ekran} & \nt{Do czego służy} & \nt{Przejścia do} \\
E1 & Start & oś nawigacji: wybór zadania & E2, E5 \\
E2 & Lista sal & lista sal z filtrem: data, godziny, liczba osób & E3, E1 \\
E3 & Szczegóły sali & plan dnia i wyposażenie sali & E4, E2 \\
E4 & Rezerwacja & formularz rezerwacji z walidacją & E6, E3 \\
E5 & Moje rezerwacje & lista własnych rezerwacji, anulowanie & E1 \\
E6 & Potwierdzenie & podsumowanie rezerwacji & E5 \\
\bottomrule
\end{tabularx}
\end{table}

\subsection{Mapa ekranów}
\zrzut{mapa_ekranow.png}{Mapa ekranów i przejść (linia ciągła: dalej, przerywana: powrót)}{0.95}

\subsection{Model nawigacji}
\textbf{Model:} oś i szprychy.\\
\textbf{Uzasadnienie:} Start to oś, a szprychy to dwa zadania: wyszukanie i rezerwacja sali oraz przegląd własnych rezerwacji. Pasuje on do użytkownika, ponieważ ma on jeden cel, więc po każdym zadaniu chce wrócić na start, a nie poznawać strukturę ekranów.

\setcounter{section}{5}
\section{Zastosowane wzorce}\label{sec:wzorce}

{\small
\begin{longtable}{@{}p{0.6cm}p{3.7cm}p{2.6cm}p{6.1cm}p{1.1cm}@{}}
\caption{Wzorce zastosowane w projekcie}\label{tab:wzorce}\\
\naglowektabeli \nt{\#} & \nt{Wzorzec} & \nt{Ekran} & \nt{Dlaczego tutaj} & \nt{Wykł.} \\
\endfirsthead
\naglowektabeli \nt{\#} & \nt{Wzorzec} & \nt{Ekran} & \nt{Dlaczego tutaj} & \nt{Wykł.} \\
\endhead
\bottomrule
\endfoot
1 & Wskazane Punkty Startowe & Start & Od razu wie, od czego zacząć & 1 \\
2 & Wybór Dwupanelowy & Lista sal na komputerze & Porównuje sale bez dodatkowych ekranów & 2 \\
3 & Uszczegółowienie W Jednym Oknie & Lista sal na telefonie & Dwa panele nie mieszczą się na wąskim ekranie & 2, 4 \\
4 & Poprawne Wartości Domyślne & Formularz i filtr & Większość rezerwacji jest na dzisiaj i na zaraz & 3 \\
5 & Zlokalizowane Komunikaty O Błędach & Formularz: przy polach godzina i liczba osób & Wiadomo, które pole i jak je poprawić & 3 \\
6 & Pionowy Stos & Filtr, lista, formularz na telefonie & Jedna kolumna & 4 \\
\end{longtable}}


\appendix
\section{Makiety z punktu kontrolnego K1}\label{zal:makiety}

\zrzut{makieta_lista_komputer.png}{Makieta z K1: lista sal (komputer)}{0.85}
\zrzut{makieta_formularz.png}{Makieta z K1: formularz rezerwacji}{0.5}
\zrzut{makieta_lista_telefon.png}{Makieta z K1: lista sal (telefon)}{0.4}


\end{document}
